% ECONOMICS_PROBLEM: {"id": "G28-235", "title": "Sovereign–bank feedback loops", "jel_code": "G28", "source_id": "OP-0255"}
% !TeX program = xelatex
\documentclass[11pt,a4paper]{article}
\usepackage[margin=23mm]{geometry}
\usepackage{fontspec}
\setmainfont{Latin Modern Roman}
\usepackage{amsmath,amssymb,mathtools}
\usepackage{xcolor,enumitem,booktabs,longtable,array,xurl,hyperref}
\definecolor{muted}{HTML}{53636C}
\hypersetup{colorlinks=true,urlcolor=blue,linkcolor=blue}
\setlength{\parindent}{0pt}
\setlength{\parskip}{5pt}
\newcommand{\R}{\mathbb R}
\newcommand{\E}{\mathbb E}
\newcommand{\N}{\mathbb N}
\newcommand{\ind}{\mathbf 1}
\newcommand{\Var}{\operatorname{Var}}
\newcommand{\Cov}{\operatorname{Cov}}
\newcommand{\supp}{\operatorname{supp}}
\DeclareMathOperator*{\argmin}{arg\,min}
\DeclareMathOperator*{\argmax}{arg\,max}
\newcommand{\entrymeta}[3]{{\sffamily\small\color{muted}\textbf{\texttt{#1}}\quad|\quad #2\par #3\par}}
\newcommand{\Problemref}[1]{\texttt{#1}}
\newcommand{\JELref}[2]{\texttt{#2} (JEL \texttt{#1})}
\begin{document}
\section*{Sovereign–bank feedback loops}
\textbf{Problem ID:} \texttt{G28-235}\quad \textbf{JEL:} \texttt{G28}\par
% BEGIN ECONOMICS STATEMENT
\label{problem:OP-0255}
{\sffamily\small\color{muted}Original subject group: Fiscal policy and sovereign debt\par}
\label{definitions:problem:30007536}
\textbf{Audit classification:} Published research agenda. Sovereign-bank linkages are a survey agenda; the prudential optimization is editorial. Verification concerns the cited version, not first priority or proof that this exact editorial specification remains unsolved on October 8, 2026.
\subsection*{Definitions and precise research target}
\paragraph{Editorial mathematical specification.} Consider a sovereign and competitive banks whose assets include loans \(L\) and domestic government bonds \(B^b\), funded by deposits and equity \(E\). Bank net worth changes with sovereign bond prices \(q^g\); loan supply depends on net worth and regulatory constraints. The government’s debt price in turn depends on fiscal capacity, output generated by lending, and possible bank support. Specify the joint default and bailout rules, household preferences, and the distribution of aggregate shocks before comparing policies.
Let a prudential policy \(a=(w,c,d)\) set the risk weight \(w\) on sovereign exposure, concentration limit \(c\), and required diversification \(d\). For feasible policies, define
\[
W(a)=\mathbb{E}\left[\sum_{t=0}^{H}\beta^t u(C_t(a),N_t(a))-\ell^{nc}(a)\right],
\]
where \(\ell^{nc}(a)\) denotes utility-equivalent nonconsumption resolution losses over the horizon, \(C_t\) is aggregate consumption after real resolution expenses, \(N_t\) labor, \(u\) household utility, and \(\beta\in(0,1)\) the discount factor. Nonconsumption resolution losses use the same utility scale and exclude expenses already deducted from consumption. Characterize policies that weaken the endogenous amplification from sovereign losses to banks and back without excessive reductions in productive credit or safe-asset services. Quantify the frontier between crisis-loss reduction and normal-time financing costs using identified exposure and lending responses. Address banks’ endogenous portfolio substitution and government incentives to provide guarantees. A policy comparison that treats sovereign bond risk and bank portfolios as fixed cannot establish its equilibrium welfare ranking.

For bank \(k\), impose \(q_t^gB^g_{kt}+q_t^L L_{kt}+R_{kt}=D_{kt}+E_{kt}\), with reserves \(R\), deposits \(D\), and market values in the same units. Define an exposure cap by \(q_t^gB^g_{kt}\le cE_{kt}\), a capital requirement by \(E_{kt}\ge\bar\kappa(wq_t^gB^g_{kt}+w_LL_{kt})\), and a diversification constraint on declared issuer portfolio shares. Each comparison recomputes sovereign prices, bank portfolios, deposits, lending, output, and feasible government guarantees. Locally, a loop may be measured by the product of the derivative of lending with respect to sovereign prices and the derivative of sovereign prices with respect to lending; global crisis loss requires the nonlinear model, not only that product. Establish nonempty feasible policy/equilibrium sets and the same crisis-resolution accounting before maximizing welfare.
\subsection*{Variants and assumption changes}
The following are \emph{editorial variants}, not additional published open conjectures.
\begin{enumerate}
\item \emph{Risk weights versus concentration.} Compare raising sovereign risk weights with imposing a direct domestic-sovereign exposure cap, keeping capital requirements and access to substitute assets fixed. Banks may respond through different portfolios and loan supply.
\item \emph{Fiscal backstop.} Permit a budget-constrained deposit guarantee or resolution fund with a fixed funding rule. Compare its ability to reduce runs against the extra sovereign liability and its effect on equilibrium bank risk taking.
\end{enumerate}
\subsection*{References and source audit}
\begin{itemize}
\item Kris James Mitchener; Christoph Trebesch. \emph{Sovereign Debt in the 21st Century}. 2021. Locator: Section5 (sovereign-bank linkages); Section10, Conclusion, printed pp.51--52. Primary text checked. \url{https://www.ifo.de/DocDL/cesifo1_wp8959.pdf}.
\end{itemize}
Sovereign-bank linkages are a survey agenda; the prudential optimization is editorial.
\begin{enumerate}\raggedright
\item\hypertarget{definitions:source:30007536:1}{} Kris James Mitchener; Christoph Trebesch. 2021. \emph{Sovereign Debt in the 21st Century}. Section5 (sovereign-bank linkages); Section10, Conclusion, printed pp.51--52.
\url{https://www.ifo.de/DocDL/cesifo1_wp8959.pdf}.
\end{enumerate}

\subsection*{Common conventions: Source mathematical conventions}

These conventions apply to each entry unless explicitly replaced.

\textbf{Models and identification.} A primitive model \(m\) specifies all probability laws, feasible choices, information, equilibrium and measurement rules used by its target. The admissible class \(\mathcal M\) is fixed independently of outcomes. If \(P_O(m)\) is its observable probability law and \(\tau(m)\) the target, its identified set at observable law \(P\) is
\[
 \mathcal I_\tau(P)=\{\tau(m):m\in\mathcal M,\ P_O(m)=P\}.
\]
Point identification means that this set is a singleton. An empty set signals incompatibility of data and assumptions. Coordinate bounds need not describe the joint set. Bounds are sharp only if every retained target is attainable or arbitrarily approachable by compatible models. Finite bounds require explicit restrictions on unobserved quantities; unrestricted confounding, hidden wealth, extrapolation or arbitrary equilibrium selection can yield uninformative sets.

\textbf{Causal interventions.} A potential outcome \(Y_i(p)\) is the outcome under a complete policy regime \(p\), including timing, financing, information and market spillovers. The target population and units are fixed before treatment. With interference, use the complete assignment vector or a justified exposure mapping; individual no-interference notation alone cannot identify an economy-wide intervention. A difference of mean potential outcomes does not require the cross-world joint distribution. Conditioning on post-treatment migration, survival, enrollment or employment changes the estimand and needs separate justification.

\textbf{Statistical statements.} An estimator, confidence set, forecast or test specifies a sampling law, model class and loss/coverage criterion. Uniform \((1-\alpha)\)-coverage means \(\inf_{m\in\mathcal M}\Pr_m\{\tau(m)\in C_n\}\geq1-\alpha\), or an explicitly asymptotic version. The whole parameter space trivially has coverage, so informativeness requires a stated width, risk or power objective. A forecasting improvement is relative to a named benchmark, information set and evaluation population.

\textbf{Equilibrium and optimization.} An equilibrium comprises feasible plans, optimality, beliefs where needed, budget constraints and all market/resource clearing. The primitives or a stated benchmark define each correspondence \(\mathcal E(p;m)\). Nonemptiness is assumed or proved, never inferred just from compact policy sets. A selected-equilibrium target uses a declared selection rule; a robust target quantifies over all permitted equilibria. Use suprema or infima unless attainment is established. An \(\varepsilon\)-optimal policy is feasible and within \(\varepsilon\) of the finite optimum value. Report an empty feasible set separately.

\textbf{Welfare and accounting.} Normative utility, interpersonal normalization, social weights, numeraire, discounting and the use of public receipts are primitives. Behavioral utility need not coincide with the welfare criterion. Household welfare includes ownership income and rents once; adding profits again to compensating variation generally double-counts them. A monetary equivalent is defined by an explicit transfer and price convention, with monotonicity and range conditions. Average effects, mechanism contributions, transfers and welfare losses are different targets. Infinite sums require convergence and dynamic feasible plans require terminal or no-Ponzi conditions.

\textbf{Source versions.} A working-paper date, revision date and journal publication year are distinguished. A DOI for a journal version is not automatically the DOI of an earlier manuscript. Locators refer to the stated version and printed pages where specified. A metadata-only check does not verify a claim about a full-text conclusion. Every entry's source subsection narrows the provenance claim accordingly.
% END ECONOMICS STATEMENT
\end{document}
